\renewcommand{\theposttitle}{The Beginner's Mind and Finding Clarity}
\renewcommand{\thepostdate}{2022-07-08}
\vspace*{\fill}
\begin{center}
\includegraphics[width=\textwidth, height=0.55\textheight, keepaspectratio]{C:/Users/James/Documents/GitHub/VisualIdeasBook/data/images/cabfc77005ddd09f821a3eeab3b3b602.png}
\end{center}
\vspace{2mm}
\begin{center}
\begin{minipage}{0.92\textwidth}
\small
As a teacher, an analogy is the most effective way to communicate foreign concepts. I came up with this one as I backed out of the driveway thinking about how to explain why my students should write to clarify their thinking about what they wanted to achieve for the year. I'm flattered to have been DM'd several times over the last couple months asking about my process of coming up with visuals. It may sound weird, but writing is an essential component of visualization. Without the clarity found through exploring an idea, a visual has no soul. That said, I still feel a deep anxiety with every visual before posting, wondering "Is This Anything?"[affiliate link] much like Jerry Seinfeld thinks before a joke.
\end{minipage}
\end{center}
\vspace*{\fill}